\chapter{Proof of the upper bound and the parameter system}\label{u:sec:proof}

We prove \cref{thm:upper,cor:upper-discrete,cor:two-sided} from
\cref{u:prop:witness,u:lem:goe-norm}. The finite-state argument of
\cite[Section~7, Lemma~8]{oai-typical-upper} converts the exclusion of
balanced cuts with small flow into a bound on the average distance from
equilibrium. We give its proof, then establish the estimate for the
discrete chain. In \cref{u:sec:budget}, we check the parameter
comparisons used throughout \cref{part:upper}.

Throughout, the disorder, Gibbs measure, and continuous-time semigroup
are those of \cref{u:sec:primitive}. In particular,
$\pi_g=\pi^{J(g)}$ and $P_t^g=P_t^{J(g)}$. The exponents are fixed before
$n$ tends to infinity.

\section{From balanced-cut flow to mixing}\label{u:sec:pruning}

The next lemma is \cite[Lemma~8]{oai-typical-upper}, specialized to the
heat-bath chain. Its proof uses the coarea argument for the reversible
Cheeger inequality of Lawler and Sokal~\cite{lawler-sokal-1988} and
entropy dissipation as in Diaconis and Saloff-Coste
\cite{diaconis-saloff-coste-1996}. We keep track of the estimates on
the Dirichlet form, since they are also used in
\cref{u:sec:proof-discrete}.

\begin{lemma}[Pruning and averaged mixing]\label{u:lem:pruning}
Fix $n\ge2$, $\beta>0$, and $g$. Let $0<\epsilon<1/3$ and $\phi>0$.
Suppose that
\[
 \mathcal Q_g(S,S^c)>\phi
 \qquad\text{whenever }\epsilon\le\pi_g(S)\le1-\epsilon.
\]
Then, for every $t>0$,
\begin{equation}\label{u:eq:averaged-mixing}
 \sum_{x\in\Sigma_n}\pi_g(x)
       \norm{P_t^g(x,\cdot)-\pi_g}_{\TV}
 \le\frac{n\sqrt{2\log2}}{\phi\sqrt t}+2\epsilon.
\end{equation}
\end{lemma}

\begin{proof}
We first remove sets that fail an internal expansion inequality. On the
remaining set, we prove a Poincar\'e inequality by integrating over level
sets. Entropy dissipation for the full chain then gives a time at which
the average distance is small; contraction gives the estimate at $t$.

In this proof, we write $\pi=\pi_g$. For distinct configurations $x,y$,
let $q(x,y)=\pi(x)p_i^{J(g)}(-x_i\mid x)$ if $y=x^{i,-x_i}$, and let
$q(x,y)=0$ otherwise. By reversibility, $q(x,y)=q(y,x)$, and
\[
 \sum_{y\ne x}q(x,y)\le n\pi(x).
\]
We use the Dirichlet form
\begin{equation}\label{u:eq:pruning-dirichlet}
 \mathfrak D(u)=\sum_{\{x,y\}}q(x,y)(u(x)-u(y))^2
              =-\langle u,\mathcal L^{J(g)}u\rangle_\pi,
\end{equation}
where each unordered pair of distinct configurations occurs once and
$\langle u,v\rangle_\pi=\sum_x\pi(x)u(x)v(x)$.

For disjoint sets $A,B\subseteq\Sigma_n$, we write
\[
 \mathcal Q_g(A,B)=\sum_{x\in A}\sum_{y\in B}q(x,y).
\]
This agrees with \eqref{u:eq:cut-flow} when $B=A^c$.

\paragraph{Removing sets.}
We start with $C=\Sigma_n$. Whenever a nonempty set $U\subset C$ satisfies
\[
 \pi(U)\le\pi(C)/2,
 \qquad \mathcal Q_g(U,C\setminus U)<\phi\pi(U),
\]
we remove $U$ from $C$. The procedure terminates after finitely many
removals. Let $A$ be the union of the sets removed up to any
stage. Every edge from $A$ to its current complement was counted when
its first endpoint was removed. Summing the inequalities at those
stages, we have, whenever $A$ is nonempty,
\[
 \mathcal Q_g(A,\Sigma_n\setminus A)<\phi\pi(A).
\]
Suppose that the removed mass reaches $\epsilon$, and consider the first
removal at which this happens. Before this removal, the removed mass
was some $a<\epsilon$, and the newly removed set has mass at most
$(1-a)/2$. After it, the removed mass is at most
$(1+a)/2<(1+\epsilon)/2<1-\epsilon$. This contradicts the hypothesis on balanced
cuts, because $\phi\pi(A)\le\phi$. We have proved that
the terminal set satisfies $\pi(C)>1-\epsilon$ and
\begin{equation}\label{u:eq:internal-expansion}
 \mathcal Q_g(U,C\setminus U)\ge\phi\pi(U)
 \quad\text{if }U\subset C\text{ and }\pi(U)\le\pi(C)/2.
\end{equation}

\paragraph{A Poincar\'e inequality on $C$.}
We claim that, for every real function $u$ on $\Sigma_n$,
\begin{equation}\label{u:eq:retained-poincare}
 \sum_{y\in C}\pi(y)(u(y)-\bar u_C)^2
 \le\frac{2n}{\phi^2}\mathfrak D(u),
 \qquad
 \bar u_C=\frac1{\pi(C)}\sum_{y\in C}\pi(y)u(y).
\end{equation}
Let $m$ be a weighted median of $u$ on $C$. Each of
$w=(u-m)_+$ and $w=(m-u)_+$ has support of mass at most $\pi(C)/2$.
Integrating \eqref{u:eq:internal-expansion} over the level sets of $w^2$,
we find
\[
 \phi\sum_C\pi w^2
 \le\sum_{\{x,y\}\subset C}q(x,y)|w(x)^2-w(y)^2|.
\]
By Cauchy--Schwarz, the right side is at most
\[
 \left(\sum_{\{x,y\}\subset C}q(x,y)(w(x)-w(y))^2\right)^{1/2}
 \left(2n\sum_C\pi w^2\right)^{1/2}.
\]
Here we used $(w(x)+w(y))^2\le2(w(x)^2+w(y)^2)$ and the bound on the
outgoing rates. Dividing and squaring, we find the desired estimate for $w$ when
$\sum_C\pi w^2>0$; when this sum vanishes, the estimate is
immediate. The sum of the energies of the two parts is at most
$\mathfrak D(u)$. Since the weighted mean minimizes the squared error,
\eqref{u:eq:retained-poincare} follows.

\paragraph{Entropy dissipation.}
For a fixed starting configuration $x$, define
$f_s^x(y)=P_s^g(x,y)/\pi(y)$. We write
\[
 \Ent_\pi(f)=\sum_y\pi(y)f(y)\log f(y),\qquad
 H(\pi)=-\sum_x\pi(x)\log\pi(x),
\]
where $\pi(f)=1$ and $0\log0=0$. Reversibility implies
\begin{align*}
 \frac{\dd}{\dd s}\Ent_\pi(f_s^x)
 &=-\sum_{\{y,z\}}q(y,z)
       (f_s^x(y)-f_s^x(z))(\log f_s^x(y)-\log f_s^x(z))\\
 &\le-4\mathfrak D(\sqrt{f_s^x}).
\end{align*}
In the last step, we used
$(a-b)(\log a-\log b)\ge4(\sqrt a-\sqrt b)^2$ for $a,b>0$.
The initial entropy is $\log(1/\pi(x))$, and the final entropy is
nonnegative. Integrating and averaging over $x\sim\pi$, we find
\begin{equation}\label{u:eq:continuous-dissipation}
 \int_0^t\sum_x\pi(x)\mathfrak D(\sqrt{f_s^x})\dd s
 \le\frac14H(\pi).
\end{equation}
To justify the calculation when the density has zero entries, we start
with $(1-\eta)f_0^x+\eta$ and let $\eta\downarrow0$. Since the state
space is finite and entropy is continuous at zero, the integrated
inequality passes to this limit.

\paragraph{Averaging the distances.}
By \eqref{u:eq:retained-poincare} and
\eqref{u:eq:continuous-dissipation}, there
exists $s\in(0,t]$ such that
\begin{equation}\label{u:eq:continuous-variance}
 \begin{split}
  \sum_x\pi(x)r_x^2&\le\frac{nH(\pi)}{2t\phi^2},\\
  r_x^2&=\sum_{y\in C}\pi(y)(\sqrt{f_s^x(y)}-a_x)^2,
  \qquad
  a_x=\frac1{\pi(C)}\sum_{y\in C}\pi(y)\sqrt{f_s^x(y)}.
 \end{split}
\end{equation}
Using Cauchy--Schwarz and the fact that $a_x^2\pi(C)\le\sum_C\pi f_s^x\le1$, we have
\[
 \sum_C\pi|f_s^x-a_x^2|\le2r_x.
\]
Further, the normalization of the density gives
$|\sum_C\pi(f_s^x-1)|\le\sum_{C^c}\pi(f_s^x+1)$.
Comparing $f_s^x$ first with $a_x^2$ and then with $1$ on $C$,
and adding the complement, we have
\begin{equation}\label{u:eq:pruning-distance}
 \sum_{\Sigma_n}\pi|f_s^x-1|
 \le4r_x+2\sum_{C^c}\pi(f_s^x+1).
\end{equation}
By stationarity, we have $\sum_x\pi(x)f_s^x(y)=1$. It follows that the
last sum in the preceding display, without its factor two, has
average $2\pi(C^c)$ over $x\sim\pi$. Dividing
\eqref{u:eq:pruning-distance} by two and using Cauchy--Schwarz in $x$,
we obtain
\[
 \sum_x\pi(x)\norm{P_s^g(x,\cdot)-\pi}_{\TV}
 \le\frac{\sqrt{2nH(\pi)}}{\phi\sqrt t}+2\epsilon.
\]
Since $|\Sigma_n|=2^n$, we have $H(\pi)\le n\log2$. Total variation
contracts under $P_{t-s}^g$ by \cref{lem:model-dynamics}, which gives
\eqref{u:eq:averaged-mixing} at time $t$.
\end{proof}

\section{Proof of the main theorem}\label{u:sec:proof-upper}

We now combine the witness estimate with \cref{u:lem:pruning}. Recall
that $\epsilon_n=n^{-\ex d}$, $\phi_n=e^{-n^{1-\ex q}}$, and
$\mathcal W_n(\epsilon_n,\phi_n)$ is the event that $\norm{W(g)}_{\mathrm{op}}
\le12$ and some $\epsilon_n$-balanced cut has flow at most $\phi_n$.

\begin{proof}[Proof of \cref{thm:upper}]
Fix $\beta>1$ and $0<\omega<1/4$, and take the affine family of
\cref{u:sec:parameters} with $\varepsilon_\omega=(1/4-\omega)/20$.
The comparisons verified in \cref{u:sec:budget} imply all hypotheses
of the preceding estimates. By \cref{u:prop:witness}, we have
\begin{equation}\label{u:eq:witness-family}
 \gamma_n(\mathcal W_n(\epsilon_n,\phi_n))
 \le C_\beta n^{-\varepsilon_\omega}+e^{-c_\beta n^{2/5}}
       +\exp\{2n^{3/4+3\varepsilon_\omega}
                    -n^{3/4+4\varepsilon_\omega}\}=o(1).
\end{equation}
Outside this event and $\{\norm{W(g)}_{\mathrm{op}}>12\}$, every
$\epsilon_n$-balanced cut has flow greater than $\phi_n$. For all
sufficiently large $n$, we have $0<\epsilon_n<1/3$. We use
\eqref{u:eq:averaged-mixing} with
\[
 \epsilon=\epsilon_n,\qquad \phi=\phi_n,
 \qquad t=e^{n^{1-\omega}}.
\]
Using Markov's inequality in the Gibbs draw and
\cref{u:lem:goe-norm}, we find
\begin{equation}\label{u:eq:joint-bad-start}
 \begin{split}
 &\Pi_n\bigl\{(g,x):
       \norm{P_{e^{n^{1-\omega}}}^g(x,\cdot)-\pi_g}_{\TV}>1/4\bigr\}\\
 &\quad\le\gamma_n(\mathcal W_n(\epsilon_n,\phi_n))+2e^{-n}
       +4n\sqrt{2\log2}\,
          e^{n^{1-\ex q}-\frac12n^{1-\omega}}+8n^{-\ex d}.
 \end{split}
\end{equation}
The right side tends to zero, since $\ex d>0$ and
$\ex q-\omega=10\varepsilon_\omega>0$. By the dictionary in
\cref{u:sec:primitive}, the $(J(g),x)$-marginal of $\Pi_n$ is
$\PP_{J,X_0\sim\pi^J}$. This proves the assertion under the joint law.
Here the Gibbs draw is held fixed when the distance is formed, and
the semigroup has rate one per site.
\end{proof}

\section{Proof of the discrete-time corollary}\label{u:sec:proof-discrete}

We use the retained set and the Poincar\'e inequality from the proof of
\cref{u:lem:pruning}. What changes is the entropy dissipation
inequality, which we prove for one step of the discrete heat-bath chain. The subsequent averaging argument
is that of \cite[Section~7, Lemma~8]{oai-typical-upper}.

\begin{proof}[Proof of \cref{cor:upper-discrete}]
Fix $\beta>1$ and $0<\omega<1/4$, and choose the same exponents as
in the proof of \cref{thm:upper}. We first fix a disorder $g$
outside the exceptional events of that proof. We abbreviate
$\pi=\pi_g$ and $K=K^{J(g)}$. For a function $u$ on $\Sigma_n$, define
\[
 Q_i u(x)=\sum_{\sigma\in\{-1,1\}}
                  p_i^{J(g)}(\sigma\mid x)u(x^{i,\sigma}).
\]
Each $Q_i$ is conditional expectation given all spins except the
$i$th, and $K=n^{-1}\sum_iQ_i$. In $L^2(\pi)$, the operators $Q_i$
are orthogonal projections. By \eqref{u:eq:pruning-dirichlet},
\begin{equation}\label{u:eq:discrete-dirichlet}
 \mathfrak D(u)=\sum_{i=1}^n\langle u,(I-Q_i)u\rangle_\pi.
\end{equation}

Let $f\ge0$ satisfy $\pi(f)=1$. The convexity of $s\log s$ implies that
$\Ent_\pi(Kf)\le n^{-1}\sum_i\Ent_\pi(Q_i f)$.
We claim that
\[
 \Ent_\pi(f)-\Ent_\pi(Q_i f)
 \ge\langle\sqrt f,(I-Q_i)\sqrt f\rangle_\pi.
\]
Indeed, on each fiber obtained by fixing all spins except the $i$th,
write $a=Q_i f$. For $s\ge0$ and $a>0$,
\[
 s\log(s/a)-s+a\ge(\sqrt s-\sqrt a)^2.
\]
This follows by writing the difference as
$2a\sqrt{s/a}(\sqrt{s/a}\log\sqrt{s/a}-\sqrt{s/a}+1)\ge0$,
with its continuous interpretation at $s=0$. Averaging on the fiber
eliminates $-s+a$. The resulting squared error is at least the squared
error about the fiber mean $Q_i\sqrt f$. Fibers with $a=0$ have $f=0$
throughout and contribute zero. Averaging over the other spins proves
the claim. Combining it with \eqref{u:eq:discrete-dirichlet}, we have
\begin{equation}\label{u:eq:discrete-entropy-step}
 \Ent_\pi(f)-\Ent_\pi(Kf)
 \ge\frac1n\mathfrak D(\sqrt f).
\end{equation}

For $f_k^x(y)=K^k(x,y)/\pi(y)$, reversibility gives
$f_{k+1}^x=Kf_k^x$. Telescoping \eqref{u:eq:discrete-entropy-step}
and averaging over $x\sim\pi$, we obtain, for every integer $m\ge1$,
\begin{equation}\label{u:eq:discrete-dissipation}
 \sum_{k=0}^{m-1}\sum_x\pi(x)\mathfrak D(\sqrt{f_k^x})
 \le nH(\pi)\le n^2\log2.
\end{equation}
Since the disorder lies outside the exceptional events, the hypothesis
on balanced cuts holds with $\phi=\phi_n$ and $\epsilon=\epsilon_n$, and
it gives the same set $C$ and inequality
\eqref{u:eq:retained-poincare}. Then there exists
$k\in\{0,\ldots,m-1\}$ such that
\[
 \sum_x\pi(x)\sum_{y\in C}\pi(y)
       (\sqrt{f_k^x(y)}-a_x)^2
 \le\frac{2n^2H(\pi)}{m\phi_n^2},\qquad
 a_x=\frac1{\pi(C)}\sum_{y\in C}\pi(y)\sqrt{f_k^x(y)}.
\]
Using \eqref{u:eq:pruning-distance}, stationarity, and
Cauchy--Schwarz, we find
\[
 \sum_x\pi(x)d^{J(g)}_{\mathrm{disc}}(x,k)
 \le\frac{2n\sqrt{2H(\pi)}}{\phi_n\sqrt m}+2\epsilon_n.
\]
Using that the distance contracts under $K^{m-k}$, we find
\begin{equation}\label{u:eq:discrete-averaged-mixing}
 \sum_x\pi_g(x)d^{J(g)}_{\mathrm{disc}}(x,m)
 \le\frac{2\sqrt{2\log2}\,n^{3/2}e^{n^{1-\ex q}}}{\sqrt m}
       +2n^{-\ex d}.
\end{equation}
Let $m=\lfloor e^{n^{1-\omega}}\rfloor$. For all sufficiently large
$n$, we have $m\ge e^{n^{1-\omega}}/2$. The first term in
\eqref{u:eq:discrete-averaged-mixing} is at most
\[
 Cn^{3/2}\exp\{n^{1-\ex q}-\tfrac12n^{1-\omega}\},
\]
which tends to zero because $\omega<\ex q$. By Markov's inequality in
the Gibbs draw, we obtain
\begin{equation}\label{u:eq:discrete-bad-start}
 \begin{split}
 &\Pi_n\{(g,x):d^{J(g)}_{\mathrm{disc}}(x,m)>1/4\}\\
 &\quad\le\gamma_n(\mathcal W_n(\epsilon_n,\phi_n))+2e^{-n}
       +Cn^{3/2}e^{n^{1-\ex q}-\frac12n^{1-\omega}}
       +8n^{-\ex d}=o(1).
 \end{split}
\end{equation}
Using the marginal identity for $\Pi_n$, we conclude that \cref{cor:upper-discrete} holds.
\end{proof}

\section{Proof of the two-sided corollary}\label{u:sec:proof-two-sided}

The two-sided statement uses the lower bound at threshold $1/4$ and
the two upper bounds just proved.

\begin{proof}[Proof of \cref{cor:two-sided}]
Fix $\beta>1$ and $0<\omega<1/4$, and let $A_\beta$ be the coefficient
in \cref{thm:barrier-730}. With $\tau=1/4$, that theorem implies
\[
 \PP_{J,X_0\sim\pi^J}
   \{d^J(X_0,e^{A_\beta n^{7/30}})>1/4\}\longrightarrow1.
\]
By the definition of the mixing time in \eqref{eq:model-mixing-time} and
\cref{lem:model-dynamics}, this event is
$\{T^J(X_0)>e^{A_\beta n^{7/30}}\}$. Intersecting it with the event
of probability tending to one in \cref{thm:upper} proves the
continuous-time assertion.

For the discrete chain, by \eqref{eq:barrier-730-discrete} at $\tau=1/4$,
we have, with probability tending to one,
\[
 d^J_{\mathrm{disc}}(X_0,\lfloor e^{A_\beta n^{7/30}}\rfloor)>1/4.
\]
Since $k\mapsto d^J_{\mathrm{disc}}(X_0,k)$ is nonincreasing, this
implies $T^J_{\mathrm{disc}}(X_0)>\lfloor e^{A_\beta n^{7/30}}\rfloor$.
On the event in \cref{cor:upper-discrete}, we also have
$T^J_{\mathrm{disc}}(X_0)\le\lfloor e^{n^{1-\omega}}\rfloor$.
The probability of their intersection tends to one, as desired.
\end{proof}

\section{The parameter system}\label{u:sec:budget}

We verify all comparisons used in the preceding chapters. Fix
$0<\omega<1/4$ and define
$\varepsilon_\omega=(1/4-\omega)/20$. Then
$0<\varepsilon_\omega<1/80$ and
$\omega=1/4-20\varepsilon_\omega$.
The affine family from \cref{u:sec:parameters} is
\begin{equation}\label{u:eq:parameter-family}
 \begin{aligned}
 \ex a&=\tfrac12,&
 \ex f&=\tfrac12,&
 \ex k_{\rm f}&=\tfrac12-2\varepsilon_\omega,\\
 \ex D&=\tfrac14-\varepsilon_\omega,&
 \ex u_{\rm E}&=\tfrac14-2\varepsilon_\omega,&
 \ex d&=\tfrac12\varepsilon_\omega,\\
 \ex e&=\tfrac14-3\varepsilon_\omega,&
 \ex b&=\tfrac12-14\varepsilon_\omega,&
 \ex z&=\tfrac14-8\varepsilon_\omega,\\
 \ex w&=\tfrac18-3\varepsilon_\omega,&
 \ex k_{\rm b}&=\tfrac14-9\varepsilon_\omega,&
 \ex s_{\rm b}&=\tfrac14-5\varepsilon_\omega,\\
 \ex s_{\rm seq}&=\tfrac14-4\varepsilon_\omega,&
 \ex r&=\tfrac14-\tfrac{15}2\varepsilon_\omega,&
 \ex u_{\rm F}&=\tfrac14-\tfrac{39}4\varepsilon_\omega,\\
 \ex q&=\tfrac14-10\varepsilon_\omega,&
 \ex z_{\rm P}&=\tfrac3{10},&
 \ex g_{\rm rot}&=\tfrac35.
 \end{aligned}
\end{equation}
An optional intermediate loss $\ex s_{\rm F}$ can be chosen to satisfy
$2\ex u_{\rm F}-\ex r<\ex s_{\rm F}<\ex q$. We leave it unspecified;
only the strict comparison $2\ex u_{\rm F}-\ex r<\ex q$ enters the
probability estimates.

The sufficient system consists of the positivity of the exponents and
\begin{equation}\label{u:eq:comparison-system}
 \begin{gathered}
  \ex a>\ex f/4,\qquad \ex a>\ex b/4,\qquad
  \ex a+\ex f/2<1,\\
  \ex k_{\rm f}<\min(\ex a,\ex f,1-\ex a),\qquad \ex b<\ex f,\\
  0<\ex k_{\rm b}<\ex z<\min(\ex a,\ex b/2,2\ex w),\\
  \max\{2\ex w,\ex w+\ex b/4,\ex b/2,
          2\ex z-\ex a,2\ex z-2\ex w,2\ex z-\ex b\}<\ex s_{\rm b},\\
  \ex a<1-\ex s_{\rm b},\qquad
  \ex e+\ex d<\ex u_{\rm E}<\min(\ex D,\ex k_{\rm f}-\ex D),\\
  \ex q<\ex u_{\rm F}<\ex r<\ex b/2,\qquad
  2\ex u_{\rm F}-\ex r<\ex q,\qquad \ex q<1/2,\\
  \ex q<\min(\ex k_{\rm b},\ex D,\ex z_{\rm P}),\qquad
  0<\ex z_{\rm P}<1/2,\\
  \ex q<\ex s_{\rm b},\qquad \ex r<\ex s_{\rm b},\qquad
  \ex a<\ex g_{\rm rot}<1-\ex s_{\rm b},\\
  \ex s_{\rm b}<\ex s_{\rm seq}<\ex e,\qquad \omega<\ex q.
 \end{gathered}
\end{equation}
All exponents are below one whenever a growing cutoff or a tail
estimate requires this. The scalar, Gram, concentration, regression,
and finite-chain arguments show that these comparisons suffice.

We list the powers used in those arguments in
\cref{u:tab:budget}. A rate is the power in a negative logarithm.
Polylogarithmic denominators in rates are retained until a strict
comparison absorbs them. Query counts and expectation bounds may
have fixed polylogarithmic numerators. The last column gives the
instance $\omega=1/5$, $\varepsilon_\omega=1/400$.

\begingroup
\small
\setlength{\tabcolsep}{4pt}
\begin{longtable}{@{}>{\raggedright\arraybackslash}p{.43\textwidth}ll@{}}
\caption{Powers in the parameter system.}\label{u:tab:budget}\\
\toprule
Quantity & Family power & $\omega=1/5$\\
\midrule
\endfirsthead
\toprule
Quantity & Family power & $\omega=1/5$\\
\midrule
\endhead
\bottomrule
\endfoot
Query count $N_{\ex f},N_{\ex b}$ & $1/2$ & $1/2$\\
Fine maximum-cell saving & $3/8$ & $3/8$\\
Coarse maximum-cell saving & $3/8+7\varepsilon_\omega/2$ & $307/800$\\
Fine $nh$, $ne^{-V_{\ex f}}$, $N_{\ex f}(1+V_{\ex f})$ & $1/2$ & $1/2$\\
Fine $\sqrt{N_{\ex f}}$, $N_{\ex f}^2V_{\ex f}R_{V_{\ex f}}/n$ & $1/4$ & $1/4$\\
Fine KL cutoff $1-\ex k_{\rm f}$ & $1/2+2\varepsilon_\omega$ & $101/200$\\
Weighted variance loss $2\ex w$ & $1/4-6\varepsilon_\omega$ & $47/200$\\
Weighted scale loss $\ex w+\ex b/4$ & $1/4-13\varepsilon_\omega/2$ & $187/800$\\
Fixed Gram loss $\ex b/2$ & $1/4-7\varepsilon_\omega$ & $93/400$\\
Weighted raw rate $1-2\ex w$ & $3/4+6\varepsilon_\omega$ & $153/200$\\
Weighted scale raw rate $1-\ex w-\ex b/4$ & $3/4+13\varepsilon_\omega/2$ & $613/800$\\
Fixed Gram raw rate $1-\ex b/2$ & $3/4+7\varepsilon_\omega$ & $307/400$\\
Joint lag / diagonal base rate $1+\ex a-2\ex z$ & $1+16\varepsilon_\omega$ & $26/25$\\
Gaussian coefficient rate $1+2\ex w-2\ex z$ & $3/4+10\varepsilon_\omega$ & $31/40$\\
Gaussian terminal rate $1+\ex b-2\ex z$ & $1+2\varepsilon_\omega$ & $201/200$\\
Coarse one-run rate $1-\ex s_{\rm b}$ & $3/4+5\varepsilon_\omega$ & $61/80$\\
Weighted quadratic energy saving $2\ex w$ & $1/4-6\varepsilon_\omega$ & $47/200$\\
Coarse log-density cost $1-\ex k_{\rm b}$ & $3/4+9\varepsilon_\omega$ & $309/400$\\
Both expected reverse-KL costs & $3/4+\varepsilon_\omega$ & $301/400$\\
Outer reverse-KL cutoff $1-\ex u_{\rm E}$ & $3/4+2\varepsilon_\omega$ & $151/200$\\
Both outer Markov decay savings & $\varepsilon_\omega$ & $1/400$\\
Raw crossing exponent $1-\ex u_{\rm E}+\ex d$ & $3/4+5\varepsilon_\omega/2$ & $121/160$\\
Crossing envelope $1-\ex e$ & $3/4+3\varepsilon_\omega$ & $303/400$\\
Total sequence failure rate $1-\ex s_{\rm seq}$ & $3/4+4\varepsilon_\omega$ & $19/25$\\
Pinning/radius power $1-\ex r$ & $3/4+15\varepsilon_\omega/2$ & $123/160$\\
Flow / good-output hit rate $1-\ex q$ & $3/4+10\varepsilon_\omega$ & $31/40$\\
Sparse-field threshold $1-\ex u_{\rm F}$ & $3/4+39\varepsilon_\omega/4$ & $1239/1600$\\
Hamming-ball entropy power $1-\ex r$ & $3/4+15\varepsilon_\omega/2$ & $123/160$\\
Bad-field rate $1+\ex r-2\ex u_{\rm F}$ & $3/4+12\varepsilon_\omega$ & $39/50$\\
Intermediate bad-field rate (optional) & see text & $(31/40,39/50)$\\
Target pressure cutoff / rate & $7/10,\ 2/5$ & $7/10,\ 2/5$\\
Regression log-cost $1-\ex D$ & $3/4+\varepsilon_\omega$ & $301/400$\\
Rotation logarithmic grid size & $3/5$ & $3/5$\\
Physical time exponent $1-\omega$ & $3/4+20\varepsilon_\omega$ & $4/5$\\
\end{longtable}
\endgroup
In the optional row, the last entry is the interval of allowed values.
In the family, that interval is
$(3/4+10\varepsilon_\omega,3/4+12\varepsilon_\omega)$.

For example, $\max\{2\ex w,\ex w+\ex b/4\}=2\ex w$, since the
difference of the two entries is $\varepsilon_\omega/2$. The weighted
variance rate and the fixed-Gram rate exceed the one-run rate by
$\varepsilon_\omega$ and $2\varepsilon_\omega$, respectively. Each of
these rates exceeds the net exponent $\ex a=1/2$ with a fixed positive
gap. We first establish the coarse event with rate $1-\ex s_{\rm b}$,
then take a union over the rotation grid. Since
$\ex s_{\rm seq}-\ex s_{\rm b}=\varepsilon_\omega$ and
$\ex e-\ex s_{\rm seq}=\varepsilon_\omega$, the grid union and the
comparison with the rare crossing each retain a strict gap. The
target pressure failure, whose rate is $2/5$, is charged once outside
that comparison.

The parameter instance at $\omega=1/5$ is
\begin{equation}\label{u:eq:parameter-instance}
 \begin{aligned}
 \omega&=\tfrac15,& \ex q&=\tfrac9{40},&
 \ex a&=\tfrac12,& \ex f&=\tfrac12,\\
 \ex k_{\rm f}&=\tfrac{99}{200},& \ex D&=\tfrac{99}{400},&
 \ex u_{\rm E}&=\tfrac{49}{200},& \ex d&=\tfrac1{800},\\
 \ex e&=\tfrac{97}{400},& \ex b&=\tfrac{93}{200},&
 \ex z&=\tfrac{23}{100},& \ex w&=\tfrac{47}{400},\\
 \ex k_{\rm b}&=\tfrac{91}{400},& \ex s_{\rm b}&=\tfrac{19}{80},&
 \ex s_{\rm seq}&=\tfrac6{25},& \ex r&=\tfrac{37}{160},\\
 \ex u_{\rm F}&=\tfrac{361}{1600},& \ex z_{\rm P}&=\tfrac3{10},&
 \ex g_{\rm rot}&=\tfrac35.
 \end{aligned}
\end{equation}
Some of the tight comparisons are listed below; the slack is the
right side minus the left side.
\begin{equation}\label{u:eq:representative-slacks}
\begin{array}{@{}lcc@{}}
\toprule
\text{Comparison}&\text{Family slack}&\omega=1/5\\
\midrule
 \ex k_{\rm f}<\ex a&2\varepsilon_\omega&1/200\\
 \ex u_{\rm E}<\ex D,\ \ex u_{\rm E}<\ex k_{\rm f}-\ex D
     &\varepsilon_\omega&1/400\\
 \ex e+\ex d<\ex u_{\rm E}&\varepsilon_\omega/2&1/800\\
 \ex z<\ex b/2&\varepsilon_\omega&1/400\\
 \ex z<2\ex w&2\varepsilon_\omega&1/200\\
 2\ex w<\ex s_{\rm b}&\varepsilon_\omega&1/400\\
 \ex w+\ex b/4<\ex s_{\rm b}&3\varepsilon_\omega/2&3/800\\
 \ex b/2<\ex s_{\rm b}&2\varepsilon_\omega&1/200\\
 2\ex z-2\ex w<\ex s_{\rm b}&5\varepsilon_\omega&1/80\\
 \ex q<\ex k_{\rm b}&\varepsilon_\omega&1/400\\
 \ex q<\ex u_{\rm F}&\varepsilon_\omega/4&1/1600\\
 \ex u_{\rm F}<\ex r&9\varepsilon_\omega/4&9/1600\\
 \ex r<\ex b/2&\varepsilon_\omega/2&1/800\\
 2\ex u_{\rm F}-\ex r<\ex q&2\varepsilon_\omega&1/200\\
 \ex s_{\rm b}<\ex s_{\rm seq},\ \ex s_{\rm seq}<\ex e
     &\varepsilon_\omega&1/400\\
 \omega<\ex q&10\varepsilon_\omega&1/40\\
\bottomrule
\end{array}
\end{equation}
The least slack in the full comparison list at this instance is
$1/1600$. The witness estimate specializes to
\begin{equation}\label{u:eq:witness-instance}
 \gamma_n(\mathcal W_n(\epsilon_n,\phi_n))
 \le C_\beta n^{-1/400}+e^{-c_\beta n^{2/5}}
       +\exp\{2n^{303/400}-n^{19/25}\}=o(1).
\end{equation}
The bound \eqref{u:eq:joint-bad-start} further adds the norm exception
$2e^{-n}$, the term
$4n\sqrt{2\log2}\exp\{n^{31/40}-n^{4/5}/2\}$, and
$8n^{-1/800}$.

\begin{proposition}[Every fixed saving below one quarter]\label{u:prop:range}
The family \eqref{u:eq:parameter-family} satisfies every sufficient
comparison in \eqref{u:eq:comparison-system} for every
$0<\varepsilon_\omega<1/80$. Hence \cref{thm:upper} holds for every
fixed $0<\omega<1/4$. These estimates do not reach $\omega=1/4$.
\end{proposition}

\begin{proof}
All parameters and all differences between the right and left sides
of \eqref{u:eq:comparison-system} are affine in $\varepsilon_\omega$,
after expanding each maximum and minimum into its constituent
comparisons. Using \eqref{u:eq:parameter-family}, we find that all endpoint values
at $0$ and $1/80$ are nonnegative and that no difference is identically
zero. Each interior value is a strict convex combination of these two
endpoint values, at least one of which is positive. Every comparison
is therefore strict. This also proves that
$\omega=1/4-20\varepsilon_\omega$ is positive; zero values at the excluded
endpoints are allowed. The optional $\ex s_{\rm F}$ can be chosen
between $2\ex u_{\rm F}-\ex r$ and $\ex q$, whose difference is
$2\varepsilon_\omega$.

In particular, $\ex q-\omega=10\varepsilon_\omega>0$.
By \cref{u:prop:witness} and \eqref{u:eq:joint-bad-start}, we conclude
that the asserted range holds. All constants
are chosen after $\omega$ and $\beta$ are fixed.

We now explain the obstruction in these estimates. Fine entropy requires
$\ex k_{\rm f}<\min(\ex a,1-\ex a)\le1/2$.
Reverse smoothing and coverage require
\[
 \ex e+\ex d<\ex u_{\rm E}
 <\min(\ex D,\ex k_{\rm f}-\ex D)
 \le\ex k_{\rm f}/2<1/4.
\]
The connecting sequence needs
$\ex q<\ex s_{\rm b}<\ex s_{\rm seq}<\ex e$, and mixing needs
$\omega<\ex q$. For any parameters satisfying these strict estimates,
we have $\omega<1/4$, independently of the affine family.
At $\varepsilon_\omega=0$, the fine, coverage, concentration, and
crossing gaps vanish. Constants and logarithms cannot be absorbed by
the arguments above. This is an obstruction to these estimates; it
does not give a lower bound or an optimality statement for the SK
dynamics.
\end{proof}

A calculation in exact rational arithmetic checks the comparisons; it
lists each affine slack and its values at $0$, $1/80$, and $1/400$.
With an optional intermediate loss chosen inside the stated interval,
the calculation checks 114 strict comparisons and 37 identities for
the powers in \cref{u:tab:budget}. The proof of \cref{u:prop:range}
does not rely on this calculation: the endpoint argument there gives
the comparisons on the whole open interval.
